%% kapa: exam=1 round=22 year=2011 period=1 form=A subject=civil_law subject_ko=민법 seq=29 number=29
%% source: Q-Net 감정평가사 기출문제 (cst00309) artlSeq=5036790 file=extracted/1008/5036790/2011년 제22회 감정평가사 1차시험 기출문제/1교시.hwp
%% conversion: kordoc 4.24.51 -> tools/gen1.py (auto); review=proofread-sonnet; flags=none
\begin{question}{29}
선의취득에 관한 설명으로 옳지 않은 것은? (다툼이 있으면 판례에 의함)
\begin{choices}
  \item 선의취득에 의하여 취득할 수 있는 동산물권은 소유권과 질권이다.
  \item 양수인이 도품이나 유실물을 공개시장에서 선의로 매수하였더라도 과실이 있는 때에는 진정한 소유자에게 대가의 변상을 청구할 수 없다.
  \item 점유보조자가 자신이 보관하고 있던 시계를 횡령하여 선의의 제3자에게 매도하고 인도한 경우, 피해자는 2년 내에 그 시계의 반환을 양수인에게 청구할 수 있다.
  \item 대리인이 본인 소유가 아닌 동산을 처분하는 경우, 선의취득이 성립할 수 있다.
  \item 목적물반환청구권의 양도에 의해서도 선의취득을 할 수 있다.
\end{choices}
\end{question}
\answer{3}
